\documentclass[12pt]{article}

\usepackage[a4paper, total={6in, 8in}]{geometry}
\usepackage{textcomp}

\begin{document}
\setlength{\parindent}{0pt}

\def \t {\textrightarrow}
\def \v {\vspace{1em}}
\def \bi {\begin{itemize}}
\def \ei {\end{itemize}}
\def \s[#1] {\section*{#1}}
\def \ss[#1] {\subsection*{#1}}
\def \sss[#1] {\subsubsection*{#1}}

\s[Ermetismo]
Ermetismo \t come ogni fenomeno si matura con una preparazione e adesione consapevole \t la pscicanalisi viene anticipata in Svevo, ma viene maturata \t anche Montale si configura nell ermetismo (piu esistenzialista) \\
L'ermetico conclamato è Quasimodo \\
Stato d'animo di desolazione totale, e introduce nell'ermetismo \t parliamo di flora e definizione di ermetismo del 1930 (decennio clu è dal 20 al 30) \\
Camillo Sbarbaro introduce questo clima, che precede la devastazione totale dei conflitti mondiali \t un'attesa che si traduce in violenza \\
L'altro volto del 900 è lo smarrimento, presagito dalla crisi dell'io con Pirandello, Tozzi \t il mondo crepuscolare e futurista sono legati in questo

\ss[Taci, anima stanca di godere - Sbarbaro]
Sbarbaro mette a fuoco malessere \t e disperazione al verso 10 \\
Tratto dalla raccolta "Pianissimo" \t carattere auto imperativo \\
A se stesso \t stanco mio cor, anche lui si rivolge all'anima \\
Si evidenzia la duplicità del godere e del soffrire, che fanno parte della vita \t nella natura è insito sempre un atomo opaco di male, e il decadentismo tratteggia visione dell essere umano che reca sempre in se una parte di male \\
Serve pero il male per conoscere il bene, e viceversa \\
Il testo palpita di emozioni e di disfattismo \t eredita dalla tradizione leopardiana e d'annunziana, in cui odire e ascoltare ricordano La pioggia del pineto (odi ascolta taci) \\
Si sente uno stato che si puo definire del "nihil" \t ne rimpianto, ne ira, ne tedio \t allora cos'è? \t è un non senso, un'assenza \\
La tradizione qua è leopardiana \t e le radici dell'espressione dell'io lirico è anche petrarchesca, che anafora insistente del non (insieme ad altre forme negative, come neppur e nessuna) \\
E espressoine che ingloba la parabola della vita: l'ira, la speranza e il tedio \\
L'ira di fronte a un evento è immediata, poi subentra la speranza di una risoluzione (che fa andare oltre l'ostacolo), ma poi la non realizzazione di una nostra progettualita si realizza nel tedio (delle giornate autunnali, caro ai crepuscolari) \\
Tedio di seneca, di leopardi \t uno spartiacque tra otto e novecento, e una cultura che aleggia dietro questi versi \t ci rivela il non senso \\
Nel sentire e nel sentore di una caduta di valore, che poi si traducono nei conflitti mondiali

\v

Taci è seguito da giaci, e il rimbalzo è tra l'anima e noi \t l'anima non si differenzia dal corpo qua \\
Rassegnazione disperata \t non dovrebbe essere disperata, quando uno si rassegna ha preso consapevolezza di cio che è successo \t irrimediabilita di una situazione porta alla rassegnazione \\
La disperazione pero è ben diverso, perche uno si dispera se ha ancora un dolore \t è la dinamica che si configura tra l'odio e l'indifferenza \t sentimento contro assenza di sentimento \\
L'odio e l'amore sono percepiti da Catullo come opposti \t sono sentimenti opposti equalmente palpitanti \t ma c'è un fuoco \\
La rassegnazione è invece assenza di speranza e di dolore \t l'annullamento, da cui si parte \t immagine di totale rinuncia \\
Dove si legge un minimo di disperazione, non c'è rassegnazione \t fa andare avanti \\
Questi termini fanno intravedere un impegno o un disimpegno politico

\v

Ultima strofa \t non ci stupiremmo se il cuore fermo \t pero respiriamo, è vegetare e non vivere \t siamo rassegnati \\
E invece camminiamo, come sonnanbuli \t evocazione di pioggia dle pineto \t senza pero la vitalita, quindi abbassamento tonale \t e non c'è simbolismo, le case sono case, donne sono donne \\
Tutto che è quella che è, non guardiamo oltre \\
Pur popolato il mondo è un deserto \t deserto dell'animo della Ginestra e nichilismo di Nietzsche \\
Anello della maglia che non tiene, cosa posso guardare allora? \t sono con gli occhi asciutti, che non hanno neanche piu le lacrime per piangere \\
Nuova epoca, e visione del quotiano della parola che affossa la malinoquenza di dannunzio (e la attraversa ) \\
Flora = definizione di ermetismo \t ? definizione di crepuscolarismo (chiedere)

\v

Ermetismo, con l'altro nome, è Carlo Bò \t letteratura come vita, il manifesto esaustivo del ruolo della figura dell'intellettuale \\
Vasetto chiuso ermeticamente \t ha la caratteristica di essere isolante, chiuso, non immediato nell'apertura \t la definizione di ermetismo come quella di una lirica chiusa, ripiegata sull'io, profondamente schiva e lontana dall'esposizione d'annunziana \\
E intrisa di una retorica che si configura in una retorica privilegiante \t che configura strane analogie \t che sono la modalita per proporre legami che nella realta non sono individuabili \\
Il poeta crea dei legami tra se e la realta, la quale non offre pero ganci al poeta \t privilegia quindi l'oscuro senso della vita, e una poesia che metafisicamente puo avere degli sbocchi in un messaggio cristiano \t cerca il senso, ma non trova il senso (se non nel non senso) \\
È una poesia che si presenta dal punto di vista esistenziale non risolutiva per l anima, ma parla del male dell'animo (il male di vivere di montale) \\
Quindi il poeta chi è? \t siamo di fronte ora a ?? \\
L ermetismo è fatto di immagini e crisi valoriale, e crisi verso la poltica (ma con testi apolitici) \t tutto cio che passa come assenza di valori nella crisi del 900 \\
I nomi che ricorrono sono 
\bi
  \item Ungaretti (avventura della parola e allegria di naufragi, poeta palombaro, scalo nell'abisso)
  \item Montale (il muro e il limite, l'anello della maglia)
  \item Salvatore Quasimodo (esprime il volto dell'ermetismo in quanto tale, che nel suo senso primario è quello di una poesia chiusa, inafferrabile e indecifrabile)
\ei
L'ermetismo è la corruzione di tutte le certezze

\v

I centri sono il bar di Firenze, e le riviste "Frontespizio", "Solaria" etc. \t molte di queste liriche escono nelle riviste \\
I poeti si mostrano nelle riviste e nei canali mediatici del tempo \t interviste televisive e radiofoniche e nei giornali \\
Vedere un autore che ci parla azzera lo scarto culturale \\
Non sta nella torre d'avorio, proprio quando montale pratica il non impegno \\

\s[Ungaretti]
Dati importanti sono la morte del figlio (come carducci), e gli studi filologici (nasce in Brasile) \t la parabola di intellettuale che vive la dimesnione dell'estraneita \\
Come Pirandello pubblica ? \t cosi Ungaretti invita a raccogliere tutte le sue liriche con le diverse versioni filologiche \t è filologo di se stesso, e porta alla poesia pura = espressione \\
Avventura della parola e poeisa pura = poesia fatta di illuminazioni dirette \t e la poesia vede scarnificare la parola, che torna alla purezza che d'annuznio aveva inquinato di sovrasensi \\
Il futurismo non è propositivo \t e in Ungaretti l'equlibrio tra spazi bianchi e parola scritta diviene espressivo \t un'immediatezza che giustifica l'espressione "poesia pura" \t la poesia è stata epurata \\
La parola non rimanda ad altro \t se non ha descrivere il paese desolato che è il suo cuore \t e l'esperienza della guerra è per una poesia che scava nell'abbisso, da cui porta l'io dell uomo \\
Non + traccia di dannunzio \t la liritica di Ungaretti attraversa e supera d'annunzio \\
La guerra è punto di riferimento \t vive la trincea come soldato, vive la morte \t scrive sui fogli che trova, con parole lapidarie che focalizzano l'urgenza del disfacimento totale in cui afferma che non è mai stato attaccato alla vita paradossalmente
\end{document}
